\chapter{2014年重庆卷（文）}

\section{单选题}

\begin{problem}
实数为一 2. 虚部为 1 的复数所对应的点位于复平面的（\quad）

\choices
    {第一象限}
    {第二象限}
    {第三象限}
    {第四象限}
\end{problem}

\begin{problem}
在等差数列 \(\{a_{n}\}\) 中，\(a_{1} = 2, a_{9} + a_{5} = 10\)，则 \(a_{7} =\)（\quad）

\choices
    {5}
    {8}
    {10}
    {14}
\end{problem}

\begin{problem}
某中学有高中生 3500 人，初中生 1500 人，为了解学生的学习情况，用分层抽样的方法从该校学生中抽取一个容量为 n 的样本，已知从高中生中抽取 70 人，则 n 为（\quad）

\choices
    {100}
    {150}
    {200}
    {250}
\end{problem}

\begin{problem}
下列函数为偶函数的是（\quad）

\choices
    {\( f(x) = x - 1 \)}
    {\(f(x) = x^{2} + x\)}
    {\( f(x)=2^{x}-2^{-x} \)}
    {\(f(x) = 2^{x} + 2^{-x}\)}
\end{problem}

\begin{problem}
执行如图所示的程序框图，则输出 s 的值为（\quad）

\choices
    {10}
    {17}
    {19}
    {36}
\end{problem}

\begin{problem}
已知命题 \( p \): 对任意 \( x \in \R \)，总有 \( |x| \geqslant 0 \)，\( q: x = 1 \) 是方程 \( x + 2 = 0 \) 的根，则下列命题为真命题的是（\quad）

\choices
    {\(p \wedge \neg q\)}
    {\(\neg p \land q\)}
    {\(\neg p \land \neg q\)}
    {\(p \wedge q\)}
\end{problem}

\begin{problem}
某几何体的三视图如图所示，则该几何体的体积为（\quad）

正视图

左视图

俯视图（\quad）

\choices
    {12}
    {18}
    {24}
    {30}
\end{problem}

\begin{problem}
设 \(F_{1}, F_{2}\) 分别为双曲线 \(\frac{x^{2}}{a^{2}} - \frac{y^{2}}{b^{2}} = 1 (a > 0, b > 0)\) 的左、右焦点，双曲线上存在一点 \(P\) 使得 \((|PF_{1}| - |PF_{2}|)^{2} = b^{2} - 3ab\) ，则该双曲线的离心率为（\quad）

\choices
    {\(\sqrt{2}\)}
    {\(\sqrt{15}\)}
    {4}
    {\(\sqrt{17}\)}
\end{problem}

\begin{problem}
若 \(\log_4(3a + 4b) = \log_4\sqrt{ab}\)，则 \(a + b\) 的最小值是（\quad）

\choices
    {\(6 + 2\sqrt{3}\)}
    {\(7 + 2\sqrt{3}\)}
    {\(6 + 4\sqrt{3}\)}
    {\(7 + 4\sqrt{3}\)}
\end{problem}

\begin{problem}
已知函数 \( f(x) = \begin{cases} \frac{1}{x + 1} - 3, & x \in (-1,0], \\ x \in (0,1], & \text{且} g(x) = f(x) - mx - m \end{cases} \) 在 \((-1,1]\) 内有且仅有两个不同的零点，则实数 \(m\) 的取值范围是（\quad）

\choices
    {\(\left(-\frac{9}{4}, -2\right] \cup \left(0, \frac{1}{2}\right]\)}
    {\(\left(-\frac{11}{4}, -2\right] \cup \left(0, \frac{1}{2}\right]\)}
    {\(\left(-\frac{9}{4}, -2\right] \cup \left(0, \frac{2}{3}\right]\)}
    {\(\left(-\frac{11}{4}, -2\right] \cup \left(0, \frac{2}{3}\right]\)}
\end{problem}

\section{填空题}

\begin{problem}
已知集合 \(A = \{3, 4, 5, 12, 13\}\)，\(B = \{2, 3, 5, 8, 13\}\)，则 \(A \cap B = \fillinblank{}\).
\end{problem}

\begin{problem}
已知向量 a 与 b 的夹角为 \( 60^{\circ} \) ，且 \( \bs{a} = (-2, -6) \) ， \( |\bs{b}| = \sqrt{10} \) ，则 \( a \cdot b = \)  \fillinblank{}.
\end{problem}

\begin{problem}
将函数 \( f(x) = \sin (\omega x + \varphi) \left( \omega > 0, -\frac{\pi}{2} \leqslant \varphi < \frac{\pi}{2} \right) \) 图象上每一点的横坐标幅度为原来的一半，纵坐标不变，再向右平移 \( \frac{\pi}{6} \) 个单位长度得到 \( y = \sin x \) 的图象，则 \( f\left(\frac{\pi}{6}\right) = \)  \fillinblank{}.
\end{problem}

\begin{problem}
已知直线 \(x - y + a = 0\) 与圆心为 \(C\) 的圆 \(x^{2} + y^{2} + 2x - 4y - 4 = 0\) 相交于 \(A, B\) 两点，且 \(AC \perp BC\) ，则实数 \(a\) 的值为 \fillinblank{} \fillinblank{}.
\end{problem}

\begin{problem}
某校早上 8:00 上课，假设该校学生小张与小王在早上 7:30\textasciitilde{}7:50 之间到校，且每人在该时间段的任何时刻到校是等可能的，则小张比小王至少早 5 分钟到校的概率为 \fillinblank{}. (用数字作答)
\end{problem}

\section{解答题}

\begin{problem}
已知 \(\{a_{n}\}\) 是首项为 1，公差为 2 的等差数列，\(S_{n}\) 表示 \(\{a_{n}\}\) 的前 \(n\) 项和. (1) 求 \(a_{n}\) 及 \(S_{n}\); (1) 求 \(a_{n}\) 及 \(S_{n}\); (2) 设 \( \{b_{n}\} \) 是首项为 2 的等比数列，公比 q 满足 \( q^{2}-(a_{1}+1)q+S_{4}=0 \) ，求 \( \{b_{n}\} \) 的通项公式及其前 n 项和 \( T_{n} \) .
\end{problem}

\begin{problem}
20 名学生某次数学考试成绩 (单位: 分) 的频率分布直方图如图.

(1) 求频率分布直方图中 a 的值; (2) 分别求出成绩落在 \( [50,60) \) 与 \( [60,70) \) 中的学生人数； (3) 从成绩在 \( [50,70) \) 的学生中任选 2 人，求此 2 人的成绩都在 \( [60,70) \) 中的概率.
\end{problem}

\begin{problem}
在 \(\triangle ABC\) 中，内角 \(A, B, C\) 所对的边分别为 \(a, b, c\)，且 \(a + b + c = 8\). (1) 若 \(a = 2, b = \frac{5}{2}\)，求 \(\cos C\) 的值; (2) 若 \(\sin A \cos^2 \frac{B}{2} + \sin B \cos^2 \frac{A}{2} = 2 \sin C\)，且 \(\triangle ABC\) 的面积 \(S = \frac{9}{2} \sin C\)，求 \(a\) 和 \(b\) 的值.
\end{problem}

\begin{problem}
已知函数 \( f(x) = \frac{x}{4} + \frac{a}{x} - \ln x - \frac{3}{2} \)，其中 \( a \in \R \)，且曲线 \( y = f(x) \) 在点 \( (1, f(1)) \) 处的切线垂直于 \( y = \frac{1}{2} x \). (1) 求 \(a\) 的值; (2) 求函数 \( f(x) \) 的单调区间和极值.
\end{problem}

\begin{problem}
如图，四棱锥 \(P - ABCD\) 中，底面是以 \(O\) 为中心的菱形，\(PO \perp\) 底面 \(ABCD, AB = 2, \angle BAD = \frac{\pi}{2}, M\) 为 \(BC\) 上一点，且 \(BM = \frac{1}{2}\). (1) 证明: \(BC \perp\) 平面 \(POM\); (2) 若 \(MP \perp AP\)，求四棱锥 \(P - ABMO\) 的体积.
\end{problem}

\begin{problem}
设椭圆 \(\frac{x^2}{a^2} + \frac{y^2}{b^2} = 1 (a > b > 0)\) 的左、右焦点分别为 \(F_1, F_2\)，点 \(D\) 在椭圆上，\(DF_1 \perp F_1F_2\)，\(\frac{|F_1F_2|}{|DF_1|} = 2\sqrt{2}, \triangle DF_1F_2\) 的面积为 \(\frac{\sqrt{2}}{2}\). (1) 求该椭圆的标准方程; (2) 是否存在圆心在 \(t\) 与 \(t\) 轴上的圆，使圆在 \(x\) 轴的上方与椭圆有两个交点，且圆在这两个交点处的两条切线相互垂直并分别过不同的焦点\(\perp\) 若存在，求出圆的方程; 若不存在，请说明理由.
\end{problem}

